\chapter{Model Background and Motivation}
\label{ch:model-background}

Yield curves are highly structured objects. Although they are observed across many maturities, much of their variation can typically be explained by only a few common drivers. This observation is central to classical approaches to term-structure modelling. Parametric models, such as those in the Nelson--Siegel family, represent the curve through a small number of interpretable parameters, while principal component analysis provides a linear decomposition of curve movements. In both cases, the underlying idea is that the term structure can be summarized by a low-dimensional representation, often associated with level, slope and curvature.

Neural-network-based approaches build on the same idea, but replace the fixed linear or parametric representation with a more flexible nonlinear one. Autoencoder models are particularly natural in this setting. An encoder maps the observed curve into a low-dimensional latent space, while a decoder reconstructs the curve from this compressed representation. Applied to yield-curve modelling, this allows the relevant factors to be learned directly from the data rather than imposed in advance.

The work of Kondratyev (2018) and Sokol (2022) shows that autoencoder neural networks can reconstruct empirically observed yield curves using only a small number of latent factors. Their results support the view that two or three factors are sufficient to describe most of the relevant variation in the curve. This is consistent with the intuition from principal component analysis, but the neural-network formulation allows the mapping from latent factors to the observed curve to be nonlinear.

A further empirical insight from this line of work is that the same network can be trained across several currencies. Instead of estimating a separate model for each currency, curves from different markets may be pooled and treated equally during training. This suggests that the underlying structure of yield curves is sufficiently similar across currencies for a common representation to be useful. Pooling currencies also increases the amount of training data, which is valuable when estimating flexible neural-network models.

A key distinction between different neural-network approaches is the object reconstructed by the decoder. In a direct autoencoder, the output can be interpreted as the curve itself. Andreasen (2023) instead connects the neural-network representation to zero-coupon bond prices, which are then converted into par swap rates using standard fixed-income pricing relations. This is important because zero-coupon bond prices are the natural pricing objects in term-structure theory. They provide a bridge between statistical curve reconstruction and the no-arbitrage restrictions required in risk-neutral modelling.

This leads to the main issue addressed by the model implemented in this thesis. Standard autoencoder models are flexible and may fit observed curves well, but they are not automatically arbitrage-free. For purely statistical curve reconstruction, this may be acceptable. However, when the model is used in a risk-neutral setting, absence of arbitrage is a structural requirement. A model may reconstruct market curves accurately and still generate term-structure dynamics that are inconsistent with the pricing restrictions of financial theory.

The affine-inspired encoder--decoder model of Matin, Wan and Poulsen addresses this limitation by combining the flexibility of autoencoders with the structure of arbitrage-free term-structure models. As in earlier neural-network approaches, the encoder maps observed par swap rates into a low-dimensional latent representation. The decoder, however, is designed differently. Instead of directly outputting swap rates, it generates zero-coupon bond prices through a construction inspired by affine term-structure models. These bond prices are then converted into par swap rates.

The central advantage of this construction is that the no-arbitrage condition is imposed as a hard constraint in the decoder. For each encoded latent state, the model solves a system of ordinary differential equations that determines the functions used in the zero-coupon bond price representation. As a result, the generated bond price curve is arbitrage-free by construction, rather than only being encouraged to satisfy no-arbitrage approximately through a penalty term or a post-processing step.

This motivates the model considered in the following chapter. Earlier neural-network approaches show that yield curves can be represented accurately using a small number of latent factors, and that such representations can be shared across currencies. However, they do not fully resolve the issue of arbitrage-free risk-neutral modelling. The affine-inspired encoder--decoder model is designed to address precisely this point: it keeps the data-driven factor representation of autoencoders, while embedding the no-arbitrage structure directly into the decoding stage.